Una empresa produeix dos tipus de peces, de ferro i d'acer. El 60\,\% de la producció total correspon a peces
de ferro i la resta són d'acer. Sabem que el 95\,\% de les peces de ferro produïdes no tenen cap defecte,
mentre que el 3\,\% de les peces d'acer són defectuoses.

\begin{apartats}

\apartat{0,75}
Si agafem una peça a l'atzar, quina és la probabilitat que sigui defectuosa?

\begin{solucio}
Considerem els successos aleatoris següents: $A$ = «la peça escollida és d'acer», $F$ = «la peça escollida és
de ferro» i $D$ = «la peça escollida és defectuosa». De les dades de l'enunciat tenim $P(F)=0{,}6$,
$P(A)=1-0{,}6=0{,}4$, $P(D\mid F)=0{,}05$ i $P(D\mid A)=0{,}03$. Per la llei de probabilitats totals, la
probabilitat que una peça sigui defectuosa és
\[
P(D)=P(D\mid F)\,P(F)+P(D\mid A)\,P(A)=0{,}05\cdot0{,}6+0{,}03\cdot0{,}4=0{,}042 .
\]

\textit{Pauta oficial:} 0,25 per plantejar la llei de la probabilitat total (o fer l'arbre de decisió o
similar); 0,25 per identificar correctament les dades del problema, i 0,25 pel càlcul concret.
\end{solucio}

\apartat{0,75}
L'empresa aviat diversificarà la producció i començarà a produir també peces de titani, que es vendran en
paquets de 5. Si la probabilitat que una peça de titani sigui defectuosa és un valor desconegut $p$, i cada
peça és defectuosa independentment de les altres, comproveu que l'expressió que ens dona la probabilitat que
en un paquet de 5 peces n'hi hagi exactament 4 de defectuoses (en funció de $p$) és $f(p)=5\left(p^4-p^5\right)$.

\begin{solucio}
Anomenem $X$ el nombre de peces defectuoses en un paquet. De l'enunciat deduïm que $X$ segueix una llei
binomial amb $n=5$ i probabilitat $p$. Per tant, la probabilitat que en un paquet hi hagi exactament quatre
peces defectuoses és
\[
P(X=4)=\binom54p^4(1-p)=\frac{5!}{4!}\,p^4(1-p)=5p^4(1-p)=5\left(p^4-p^5\right).
\]
També podem argumentar directament, sense fer servir la llei binomial. Com que hi ha quatre peces
defectuoses, n'hi haurà exactament una que no ho és. La probabilitat que una en concret (per exemple, la
primera) sigui bona i les altres quatre defectuoses és $(1-p)p^4$. Com que hi ha cinc possibilitats per
escollir quina és la peça bona, $P(X=4)=5(1-p)p^4=5\left(p^4-p^5\right)$.

\textit{Pauta oficial:} 0,5 pel plantejament i 0,25 pel càlcul.
\end{solucio}

\apartat{1}
Considereu la funció $f(p)$ de l'apartat anterior. Determineu el valor màxim que pren $f(p)$ quan $p\ge0$.

\begin{solucio}
Busquem els extrems relatius de la funció $f(p)=5\left(p^4-p^5\right)$:
$f'(p)=5\left(4p^3-5p^4\right)=0\rightarrow p=0,\ p=\frac45$. Observem que $f(0)=0$ i que
$f\!\left(\frac45\right)>0$. Mirem el signe de la segona derivada:
\[
f''(p)=5\left(12p^2-20p^3\right)\;\rightarrow\;f''\!\left(\tfrac45\right)
=5\left(12\left(\tfrac45\right)^2-20\left(\tfrac45\right)^3\right)<0,
\]
i deduïm que a $p=\frac45$ la funció $f(p)$ assoleix un màxim local. Com que per a $p\in\left(0,\frac45\right)$
la funció creix i per a $p>\frac45$ la funció decreix, es tracta d'un màxim absolut quan $p$ és no negativa.
I el valor màxim serà $f\!\left(\frac45\right)=5\left(\left(\frac45\right)^4-\left(\frac45\right)^5\right)=0{,}41$.

\textit{Pauta oficial:} 0,25 pel càlcul de la derivada i els punts crítics, 0,5 per argumentar que el màxim
absolut en $p>0$ es troba a $p=\frac45$, i 0,25 per trobar el valor d'aquest màxim.
\end{solucio}

\end{apartats}
